% ProveIt Contributors. Classical foundations are attributed below.
\section{Bilinear moments of balanced Stieltjes kernels}\label{stbil:sec}

Translation has supplied linear identities and their boundary values.
To obtain exact product identities, combine the Stieltjes indices so that
their common logarithmic tails cancel. The balanced kernels of the
preserved \emph{Balanced Stieltjes Moments} report~\cite{stmell:balanced} are
\begin{equation}\label{stbil:kernel}
 C_p(x,a)=\left.\partial_z^p\left(x^z\zeta(1+z,x+a)-\frac1z\right)
 \right|_{z=0},\qquad x,a>0,\quad p\in\mathbb Z_{\geq0}.
\end{equation}
With $\gamma_j$ denoting the generalized Stieltjes functions and
$L=\log x$, this is the explicit combination
\[
 C_p(x,a)=\frac{L^{p+1}}{p+1}
       +\sum_{j=0}^p(-1)^j\binom pj\gamma_j(x+a)L^{p-j}.
\]
The subtraction removes the pole normally in $z$.  We first recall and
prove the analytic input from that report, rather than claiming it anew.

\begin{lemma}[Balanced scalar transform and decay]\label{stbil:scalar}
At zero $C_p(x,a)=O(1+|\log x|^{p+1})$, and at infinity
$C_p(x,a)=O(x^{-p-1})$.  The scalar Mellin transform is
\begin{align*}
 J_p(s,a)&:=\int_0^\infty x^{s-1}C_p(x,a)\,dx
       =\pi\csc(\pi s)\zeta(1-s,a)Q_p(s),
       &&0<\Re s<p+1,\\
 \sum_{p\geq0}Q_p(s)\frac{z^p}{p!}
       &=\frac{\Gamma(s+z)}{\Gamma(s)\Gamma(1+z)}.
\end{align*}
Apparent singularities at interior integers are removable in the whole
product defining $J_p$.
\end{lemma}

\begin{proof}
The explicit finite Stieltjes combination and regularity of $\gamma_j$
on positive compact intervals give the zero bound.  Euler--Maclaurin
for Hurwitz zeta, expanded about $x$ with shift $a$, gives for every $K$
\[
 x^z\zeta(1+z,x+a)-\frac1z
 =\sum_{k=1}^{K}\frac{(-1)^kB_k(a)}{k!}(1+z)_{k-1}x^{-k}
        +O(x^{-K-1}).
\]
Here $B_k$ is the Bernoulli polynomial with $B_1(a)=a-\tfrac12$.
Taking sufficiently many Euler--Maclaurin terms bounds the remainder
uniformly on a slightly larger $z$-disc; Cauchy differentiation thus
preserves its inverse-power bound.  The polynomial $(1+z)_{k-1}$ has
degree $k-1$, so its $p$th derivative vanishes for $k\leq p$.
This proves the claimed infinity bound and absolute convergence of $J_p$
on $0<\Re s<p+1$.

For the transform, the initially convergent beta identity is
\[
 \int_0^\infty x^{\alpha-1}\zeta(\beta,x+a)\,dx
 =\frac{\Gamma(\alpha)\Gamma(\beta-\alpha)}{\Gamma(\beta)}
       \zeta(\beta-\alpha,a),
 \quad\Re\alpha>0,\quad\Re\beta>\Re\alpha+1.
\]
It follows by integrating the absolutely convergent Hurwitz series.
Continue its left side one strip farther by the explicit cutoff
\begin{align*}
 &\int_0^1x^{\alpha-1}\zeta(\beta,x+a)\,dx\\
 &\quad+\int_1^\infty x^{\alpha-1}
       \left(\zeta(\beta,x+a)-\frac{x^{1-\beta}}{\beta-1}\right)\,dx
       +\frac1{(\beta-1)(\beta-\alpha-1)}.
\end{align*}
For $\Re(\beta-1)<\Re\alpha<\Re\beta$, $\Re\alpha>0$, the last term
equals $-\int_0^1x^{\alpha-\beta}\,dx/(\beta-1)$.  Thus the continued
expression is the ordinary integral of the subtracted integrand over
the whole positive axis.  Substitute $\alpha=s+z$, $\beta=1+z$ with
$0<\Re s<1$, initially $z\ne0$ sufficiently small.  This gives
\[
 \int_0^\infty x^{s-1}
       \left(x^z\zeta(1+z,x+a)-\frac1z\right)\,dx
 =\frac{\Gamma(s+z)\Gamma(1-s)}{\Gamma(1+z)}\zeta(1-s,a).
\]
Both sides extend holomorphically through $z=0$.  Near zero, the removable
term $(x^z-1)/z$ is bounded by a constant times
$(1+|\log x|)x^{-|\Re z|}$; choose $|\Re z|<\Re s/2$.  Near infinity,
the normally uniform expansion above applies.  Dominated parameter
differentiation therefore gives the claimed scalar formula on
$0<\Re s<1$, using the gamma reflection identity.
The differentiated integral is holomorphic on $0<\Re s<p+1$ by the
endpoint bounds.  Analytic continuation identifies it with the displayed
meromorphic expression throughout this larger strip, forcing its interior
poles to be removable.  No unsubtracted divergent monomial integral is
assigned a value in this argument.
\end{proof}

In particular,
\[
 Q_0=1,\quad Q_1=\psi(s)+\gamma,\quad
 Q_2=(\psi(s)+\gamma)^2+\psi'(s)-\zeta(2).
\]

\subsection{An explicit single-contour product formula}

\begin{theorem}[Bilinear Mellin--Barnes coordinates]\label{stbil:barnes}
For $a,b>0$, $p,q\geq0$, and $0<\Re v<p+q+2$, the ordinary integral
\[
 I_{p,q}(v;a,b)=\int_0^\infty x^{v-1}C_p(x,a)C_q(x,b)\,dx
\]
is absolutely convergent and holomorphic in $v$.  Choose any real $c$ with
\[
 \max(0,\Re v-q-1)<c<\min(p+1,\Re v).
\]
Then
\begin{align*}
 I_{p,q}(v;a,b)
 &=\frac1{2\pi i}\int_{c-i\infty}^{c+i\infty}
       J_p(s,a)J_q(v-s,b)\,ds\\
 &=\frac{\pi^2}{2\pi i}\int_{c-i\infty}^{c+i\infty}
 \frac{\zeta(1-s,a)\zeta(1-v+s,b)Q_p(s)Q_q(v-s)}
      {\sin(\pi s)\sin(\pi(v-s))}\,ds.
\end{align*}
The contour is absolutely and exponentially convergent.  Its value is
independent of the permissible split $c$.
\end{theorem}

\begin{proof}
The endpoint bounds give absolute convergence, locally uniformly in the
stated $v$ strip; the same argument allows arbitrary fixed powers of
$\log x$.  The interval for $c$ is nonempty precisely under the stated
strict inequalities on $\Re v$.  Along a fixed vertical line with positive
real part, the gamma-ratio derivatives $Q_p$ have at most powers of
$\log(2+|\Im s|)$ as growth.  The Hurwitz zeta function with fixed positive
shift has at most polynomial growth on every fixed vertical strip.  For
example this follows directly by choosing the Euler--Maclaurin cutoff
proportional to $1+|\Im s|$ and taking a fixed sufficiently long remainder.
The sine factor therefore implies
\[
 J_p(c+it,a)=O\bigl(e^{-\pi|t|}(1+|t|)^A
                         (1+\log(2+|t|))^p\bigr)
\]
for some $A$, and similarly for the second factor.  Any interior integer
at finite height is interpreted through the holomorphic scalar transform.

Mellin inversion of $C_p$ on this line is legitimate: in logarithmic
coordinates its transform is the Fourier transform of the continuous
integrable function $e^{cy}C_p(e^y,a)$, and $J_p(c+it,a)$ is integrable in
$t$.  Substitute that absolutely convergent inversion formula into
$I_{p,q}$.  Fubini follows from the explicit product bound
\[
 \int_{-\infty}^{\infty}|J_p(c+it,a)|\,dt\;
 \int_0^\infty x^{\Re v-c-1}|C_q(x,b)|\,dx<\infty.
\]
The inner $x$ integral is then $J_q(v-s,b)$, giving the first equality.
The second equality uses the already proved scalar formula.  It is an
application of classical Mellin inversion and its Parseval formula;
it is not obtained by multiplying two original $x$ integrals.
\end{proof}

This is a reduction to one explicit Hurwitz--gamma Barnes contour.
It does not assert a finite reduction of every product to single Hurwitz
values, classical polylogarithms, or multiple zeta values.  It does,
however, supply a reproducible exponentially convergent coordinate for
the whole ordinary product strip, including $\Re v\geq2$ when $p+q>0$.
Differentiating in $v$ gives logarithmically weighted products, with
normal convergence on compact substrips.

For real $0<v<2p+2$ there is the positive diagonal form
\[
 I_{p,p}(v;a,a)=\frac1\pi\int_0^\infty
        |J_p(v/2+it,a)|^2\,dt.
\]
At $p=0,a=1,v=1$ this specializes to the classical-type Parseval identity
\[
 \int_0^\infty(\log x-\psi(x+1))^2\,dx
 =\pi\int_0^\infty\frac{|\zeta(\tfrac12+it)|^2}
                              {\cosh^2(\pi t)}\,dt.
\]
No novelty claim is made for this scalar specialization.

\subsection{Exact moments of adjacent-shift differences}

Define
\[
 D_p(x,a)=C_p(x,a)-C_p(x,a+1).
\]
The Hurwitz recurrence gives the elementary identity, also recorded in
the balanced-Stieltjes source,
\begin{equation}\label{stbil:D}
 D_p(x,a)=\frac{\log^p(x/(x+a))}{x+a}.
\end{equation}
This exposes an exactly solvable family of linear combinations of four
balanced product moments.  These differences retain the $x^{-p-1}$
decay and have a fixed sign $(-1)^p$.

\begin{theorem}[Factorial and integer-weight bilinear identities]\label{stbil:integer-moments}
Let $p,q\geq0$, $N=p+q$, and $a>0$.  Then
\begin{equation}\label{stbil:factorial}
 \int_0^\infty D_p(x,a)D_q(x,a)\,dx
       =\frac{(-1)^N N!}{a}.
\end{equation}
For integers $2\leq m\leq N+1$, write
\[
 R_m(T)=\frac{\prod_{r=2}^{m-1}(T-r)}{(m-2)!}
       =\sum_{j=0}^{m-2}r_{m,j}T^j
\]
(the empty product for $m=2$ is $1$).  Then the ordinary convergent moment is
\begin{equation}\label{stbil:integer}
 \int_0^\infty x^{m-1}D_p(x,a)D_q(x,a)\,dx
 =(-1)^N N!a^{m-2}
       \sum_{j=0}^{m-2}r_{m,j}\bigl(\zeta(N+1-j)-1\bigr).
\end{equation}
All zeta arguments on the right are integers at least two.
\end{theorem}

\begin{proof}
Set $t=\log(1+a/x)$.  Reversing endpoints gives, for the entire convergence
strip $0<\Re v<N+2$,
\begin{equation}\label{stbil:Laplace}
 \int_0^\infty x^{v-1}D_p(x,a)D_q(x,a)\,dx
 =(-1)^N a^{v-2}\int_0^\infty
          t^Ne^{-t}(e^t-1)^{1-v}\,dt.
\end{equation}
At $v=1$ the gamma integral proves the factorial formula.  At $v=m\geq2$,
\[
 e^{-t}(e^t-1)^{1-m}
  =\sum_{\ell=m}^{\infty}\binom{\ell-2}{m-2}e^{-\ell t}.
\]
These terms are nonnegative.  Tonelli and the gamma integral therefore
give
\[
 (-1)^N N!a^{m-2}\sum_{\ell=m}^{\infty}
          \binom{\ell-2}{m-2}\ell^{-N-1}.
\]
The polynomial $R_m(\ell)$ equals the binomial coefficient here and
vanishes at $\ell=2,\ldots,m-1$.  Thus the sum may start at $\ell=2$.
Expanding the finite polynomial gives the claimed formula.  Its convergence
uses $N\geq m-1$; every resulting zeta series is absolutely convergent.
\end{proof}

For example, whenever the indicated $N$ is large enough,
\begin{align*}
 \int_0^\infty xD_pD_q\,dx
   &=(-1)^NN!\bigl(\zeta(N+1)-1\bigr), &&N\geq1,\\
 \int_0^\infty x^2D_pD_q\,dx
   &=(-1)^NN!a\bigl(\zeta(N)-2\zeta(N+1)+1\bigr), &&N\geq2,\\
 \int_0^\infty x^3D_pD_q\,dx
   &=\frac{(-1)^NN!a^2}{2}
        \bigl(\zeta(N-1)-5\zeta(N)+6\zeta(N+1)-2\bigr), &&N\geq3.
\end{align*}
Here both factors in each integrand have the same shift $a$.
In particular
\[
 \int_0^\infty x\,[C_0(x,a)-C_0(x,a+1)]
                       [C_1(x,a)-C_1(x,a+1)]\,dx=1-\zeta(2),
\]
and
\[
 \int_0^\infty x^2[C_1(x,a)-C_1(x,a+1)]^2\,dx
      =2a\bigl(\zeta(2)-2\zeta(3)+1\bigr).
\]

\subsection{The factorial Gram matrix and Laguerre coordinates}

The Gram matrix of $D_0,\ldots,D_M$ in $L^2(0,\infty;dx)$ is
\[
 G_{p,q}=\frac{(-1)^{p+q}(p+q)!}{a},\qquad
 \det G=a^{-(M+1)}\prod_{j=0}^{M}(j!)^2.
\]
Indeed the classical Laguerre polynomials
$L_n(t)=\sum_{p=0}^n(-1)^p\binom np t^p/p!$ give the triangular combinations
\[
 \Lambda_n(x,a)=\sum_{p=0}^n\binom np\frac{D_p(x,a)}{p!}
       =\frac{L_n(\log(1+a/x))}{x+a}.
\]
The same substitution as above and classical Laguerre orthogonality yield
\[
 \int_0^\infty\Lambda_m(x,a)\Lambda_n(x,a)\,dx
          =\frac{\delta_{mn}}a.
\]
The diagonal entries of the triangular change of basis are $1/n!$;
taking determinants proves the Gram determinant.  This use of Laguerre
orthogonality is classical; the balanced-kernel difference supplies the
coordinate map.

\subsection{Different shifts and an explicit hypergeometric coordinate}

For $a,b>0$ define, initially when $0<\Re v<2$ and $u,w$ are near zero,
\[
 K(v;u,w;a,b)=\int_0^\infty
       \frac{x^{v+u+w-1}}{(x+a)^{1+u}(x+b)^{1+w}}\,dx.
\]
The Euler substitution $x=at/(1-t)$ gives
\begin{equation}\label{stbil:hypergeometric}
 K(v;u,w;a,b)=
 \frac{a^{v+w-1}}{b^{1+w}}B(v+u+w,2-v)
 {}_2F_1(1+w,v+u+w;2+u+w;1-a/b).
\end{equation}
The argument $1-a/b$ is real and less than one, so the Euler integral
fixes the real positive branch unambiguously.  Differentiation yields
\[
 \int_0^\infty x^{v-1}D_p(x,a)D_q(x,b)\,dx
       =\left.\partial_u^p\partial_w^qK(v;u,w;a,b)\right|_{u=w=0}.
\]
The left side is holomorphic on $0<\Re v<p+q+2$.  Consequently the whole
parameter derivative on the right continues to this larger strip, with
all apparent interior poles removable.  The derivative is taken first
in the initial strip; one must not differentiate an undifferentiated
integral at $\Re v\geq2$ where it diverges.  This supplies a finite
hypergeometric-parameter-derivative function class for unequal shifts.

At $a=b$ it simplifies to
\[
 \int_0^\infty x^{v-1}D_p(x,a)D_q(x,a)\,dx
   =a^{v-2}\left.\partial_z^{p+q}B(v+z,2-v)\right|_{z=0},
\]
with the same continuation convention.  At $p=q=0,v=1$ and $a\ne b$,
the unequal-shift integral is $\log(b/a)/(b-a)$, with limit $1/a$.

\subsection{Attribution and evidence boundary}

The scalar $J_p$ theorem, the decay of $C_p$, and its adjacent-shift
recurrence come from the preserved
balanced-Stieltjes report, whose further-research question 5 explicitly
asks for a bilinear theory.  The single-contour formula is the direct
application of classical Mellin Parseval with the full admissible split
and the enlarged balanced convergence strip proved above.  The exact
adjacent-shift formulas use the classical Hurwitz recurrence, beta/gamma
integrals, and Laguerre orthogonality.  These are transparent derived
identities in the balanced-kernel setting, not a claim that Parseval,
Laguerre theory, or their scalar special cases are new.

Relevant classical references are NIST DLMF~\cite{stbil:classical}
\S2.5(i), especially
equations 2.5.2--2.5.5 (Mellin inversion and Parseval), \S15.6.1 (Euler's
hypergeometric integral), and Table 18.3.1 (Laguerre orthogonality).
Targeted searches of the maintained chapters and the Stieltjes-correlation
report tree found no existing balanced-product/Laguerre formulation;
this is only a corpus-overlap check, not a proof of literature novelty.

The accompanying scripts separate mathematical proofs from diagnostics.
The Barnes checks compare direct positive-axis quadrature with independent
vertical-contour quadrature.  Higher $C_p$ values use the preserved
Euler--Maclaurin evaluator, and the contour is truncated at height 24.
These are noninterval numerical diagnostics; zero printed residual does
not mean an exact or rigorously bounded zero.  Finite Gram determinants
and the integer-weight coefficient checks use exact rational arithmetic.
